\chapter*{卷首：把三国看作三种抵达方式}
\addcontentsline{toc}{chapter}{卷首：把三国看作三种抵达方式}

三国并不是三种颜色填满同一张地图。魏在洛阳面对的是怎样把华北仓储、
尚书文书和禁军送到淮南、关中与辽东；蜀汉在成都面对的是怎样把盆地、
南中与汉中接成能供养一支边军的腹地；孙吴在建业面对的是怎样把江面、
州郡和家庭簿籍转化为既能防守又能远征的能力。它们都有皇帝、诏令和
将军，却没有谁能免于道路、水面、山口、继承与地方社会的限制。

本卷以 220 年汉魏交接为入口，以 280 年西晋接收孙吴为收束。它不把
《三国演义》的后世人物关系当作史实，也不把正史的一条胜负判语变成
完整战场。读者将先进入三座中枢，再随夷陵、祁山、辽东、淮南和长江的
行动观察资源如何变成选择；最后回到走马楼吴简、石经残石与律令，追问
国家命令在何处真正碰到家庭、书写和地方惯例。

每一处事件链接都通向首次展开的节点；未列入索引的账本条目仍在
`data/event-ledger/three-kingdoms-220-280.csv` 中保存。材料有时只能
确认结果而不能确认动机，本书宁愿让这种空白留在页上，也不以一个流畅
故事替它填满。
